\chapter{Training Data Generation Engine}
\label{chap:generator}\label{sim:chap:generator}

\section{Algorithm 2: Agentic Data Synthesis}

Given taxonomies $\mathcal{T}$ and a target size $N$, the generator produces a synthetic
dataset optimized for diversity, complexity, and quality.

\mermaidfig[0.55\textwidth]{diagrams/generator-stages}%
  {Data generation stages implementing Algorithm 2.}%
  {fig:generator-stages}

\section{Step 1: Mix Sampling (Global Diversity)}

Nodes are sampled from taxonomies to create ``mixes''---combinations of factors:

\begin{lstlisting}[language=jsonx,caption={Mix sampling (simula\_data\_generator.py, lines 120--150)}]
def _sample_mix(
    self, taxonomies: list[Taxonomy], strategy: SamplingStrategy
) -> list[TaxonomyNode]:
    """Sample a mix of nodes from taxonomies."""
    mix = []
    for taxonomy in taxonomies:
        if strategy == SamplingStrategy.UNIFORM:
            # Sample uniformly from leaf nodes
            leaves = self._get_leaf_nodes(taxonomy.root)
            node = random.choice(leaves)
        elif strategy == SamplingStrategy.COVERAGE:
            # Prioritize under-sampled nodes
            node = self._get_least_sampled_node(taxonomy)
        mix.append(node)
    return mix
\end{lstlisting}

\section{Step 2: Meta-Prompt Generation (Local Diversity)}

Each mix is converted to a meta-prompt that guides example generation:

\begin{lstlisting}[language=jsonx,caption={Meta-prompt generation}]
async def _generate_meta_prompt(
    self, mix: list[TaxonomyNode], schema_context: str
) -> str:
    """Generate a meta-prompt from a taxonomy mix."""
    prompt = f"""
    Generate a natural language question that a user might ask,
    given these requirements:
    
    {self._format_mix(mix)}
    
    Schema context:
    {schema_context}
    
    The question should be realistic and answerable with SQL.
    """
    response = await self.llm.generate(prompt, temperature=0.9)
    return response.content
\end{lstlisting}

\section{Step 3: Complexification}

A fraction $c$ (default 0.5) of examples undergo complexification:

\begin{lstlisting}[language=jsonx,caption={Complexification prompt}]
SYSTEM: You increase the complexity of a Text-to-SQL example
        while maintaining correctness.
        
USER:   Original question: {question}
        Original SQL: {sql}
        
        Add complexity by:
        - Adding edge cases or constraints
        - Requiring additional joins
        - Adding subqueries or CTEs
        - Handling NULL values explicitly
        
        Output the complexified question and SQL.
\end{lstlisting}

\section{Step 4: Example Generation}

The meta-prompt is used to generate a complete training example:

\begin{lstlisting}[language=jsonx,caption={Example generation}]
async def _generate_example(
    self, meta_prompt: str, schema_context: str
) -> TrainingExample:
    """Generate a single training example."""
    prompt = f"""
    Generate a Text-to-SQL training example.
    
    Meta-prompt: {meta_prompt}
    Schema: {schema_context}
    
    Output JSON:
    {{
      "question": "...",
      "sql": "SELECT ..."
    }}
    """
    response = await self.llm.generate(prompt)
    return self._parse_example(response.content)
\end{lstlisting}

\section{Step 5: Double-Critic Filtering}

The double-critic mitigates sycophancy bias by independently querying correctness
and incorrectness:

\begin{lstlisting}[language=jsonx,caption={Double-critic evaluation (simula\_llm\_client.py, lines 180--220)}]
async def double_critic_evaluate(
    self, question: str, sql: str, schema: str
) -> tuple[bool, str]:
    """Evaluate with double-critic to mitigate sycophancy."""
    
    # Query 1: Is this correct?
    correct_response = await self.generate(
        f"Is this SQL correct for the question?\n"
        f"Question: {question}\n"
        f"SQL: {sql}\n"
        f"Schema: {schema}\n"
        f"Answer YES or NO with explanation."
    )
    
    # Query 2: Is this incorrect?
    incorrect_response = await self.generate(
        f"Is this SQL INCORRECT for the question?\n"
        f"Question: {question}\n"
        f"SQL: {sql}\n"
        f"Schema: {schema}\n"
        f"Answer YES or NO with explanation."
    )
    
    # Accept only if: correct=YES AND incorrect=NO
    is_valid = (
        "YES" in correct_response.content.upper() and
        "NO" in incorrect_response.content.upper()
    )
    
    return is_valid, f"{correct_response.content}\n---\n{incorrect_response.content}"
\end{lstlisting}

\section{Schema Context Pruning}
\label{sec:schema-context-pruning}

For production deployments with wide tables (100+ columns), the full schema context
may exceed the LLM's context window. The schema context pruning strategy selects
relevant columns to keep the prompt within token limits while preserving generation quality.

\subsection{Token Budget}

Given a target token limit $T$ (default: 4096 for schema context), allocate:
\begin{itemize}
  \item 60\% for directly relevant columns (based on taxonomy path)
  \item 25\% for structurally related columns (foreign keys, common joins)
  \item 15\% for metadata and descriptions
\end{itemize}

\subsection{Pruning Algorithm}

\begin{lstlisting}[language=jsonx,caption={Schema context pruning algorithm}]
def prune_schema_context(
    schema: SchemaRegistry,
    taxonomy_path: list[str],
    max_tokens: int = 4096,
    estimator: TokenEstimator = default_estimator
) -> str:
    """
    Prune schema context to fit within token budget.
    
    Strategy:
    1. Identify primary tables from taxonomy path keywords
    2. Include all columns from primary tables (up to 30 per table)
    3. Include foreign key relationships
    4. Include frequently-joined tables (from usage statistics)
    5. Truncate descriptions to first sentence
    """
    # Extract keywords from taxonomy path
    keywords = extract_keywords(taxonomy_path)
    
    # Score tables by relevance
    scored_tables = []
    for table in schema.tables:
        relevance = compute_relevance(table, keywords)
        scored_tables.append((table, relevance))
    
    # Sort by relevance and select top tables
    scored_tables.sort(key=lambda x: x[1], reverse=True)
    
    # Build context incrementally within budget
    context_parts = []
    token_count = 0
    
    for table, relevance in scored_tables:
        table_context = format_table_context(
            table,
            max_columns=30,
            truncate_descriptions=True
        )
        table_tokens = estimator.count(table_context)
        
        if token_count + table_tokens <= max_tokens:
            context_parts.append(table_context)
            token_count += table_tokens
        else:
            break
    
    return "\n\n".join(context_parts)
\end{lstlisting}

\subsection{Relevance Scoring}

Column relevance is computed based on:
\begin{enumerate}
  \item \textbf{Name match} (weight: 0.4): Column name contains taxonomy keywords
  \item \textbf{Type affinity} (weight: 0.2): Column type matches expected query patterns
        (e.g., DATE for temporal queries, DECIMAL for aggregations)
  \item \textbf{Structural importance} (weight: 0.2): Primary keys, foreign keys, indexes
  \item \textbf{Usage frequency} (weight: 0.2): Historical query patterns if available
\end{enumerate}

\subsection{Configuration}

\begin{table}[H]
\centering
\caption{Schema Pruning Configuration}
\label{tab:schema-pruning-config}
\begin{tabular}{lll}
\toprule
\textbf{Parameter} & \textbf{Default} & \textbf{Description} \\
\midrule
\texttt{SIMULA\_SCHEMA\_MAX\_TOKENS} & 4096 & Max tokens for schema context \\
\texttt{SIMULA\_MAX\_COLUMNS\_PER\_TABLE} & 30 & Column limit per table \\
\texttt{SIMULA\_INCLUDE\_FK\_DEPTH} & 2 & Foreign key traversal depth \\
\texttt{SIMULA\_TRUNCATE\_DESCRIPTIONS} & true & Truncate to first sentence \\
\bottomrule
\end{tabular}
\end{table}

\subsection{Quality Impact}

Schema pruning has measurable impact on generation quality:
\begin{itemize}
  \item \textbf{Aggressive pruning} ($<$2048 tokens): May miss relevant columns,
        leading to incomplete SQL. Use only for narrow-domain generation.
  \item \textbf{Standard pruning} (4096 tokens): Suitable for most production
        scenarios with tables up to 200 columns.
  \item \textbf{Full context} (no pruning): Preferred when context window permits;
        required for complex multi-table queries.
\end{itemize}

\noindent
The v1.3 roadmap includes adaptive pruning based on taxonomy complexity---complex
taxonomy paths will automatically receive larger schema budgets.

\section{Output Format}

Generated examples are written to JSONL:

\begin{lstlisting}[language=jsonx,caption={JSONL output format}]
{"id": "sim-001", "question": "What is the total revenue...", 
 "sql": "SELECT SUM(amount) FROM...", "complexity_score": 0.3, ...}
{"id": "sim-002", "question": "Show all customers with...", 
 "sql": "SELECT * FROM customers WHERE...", "complexity_score": 0.7, ...}
\end{lstlisting}

% =============================================================================
% CHAPTER 6: ENVIRONMENT AND CLI
% =============================================================================
\newpage